% ============================================================
%  Context-Aware Vocabulary Pruning for LLM Inference
%  Two-column arXiv-style  |  pdflatex
% ============================================================
\documentclass[twocolumn,10pt]{article}

% ── Geometry ────────────────────────────────────────────────
\usepackage[
  a4paper,
  top=19mm, bottom=23mm,
  left=14mm, right=14mm,
  columnsep=6mm
]{geometry}

% ── Core packages ───────────────────────────────────────────
\usepackage{amsmath,amssymb}
\usepackage{graphicx}
\usepackage{booktabs}
\usepackage{microtype}
\usepackage{xcolor}
\usepackage{url}
\usepackage{enumitem}
\usepackage{caption}
\usepackage{subcaption}
\usepackage{float}
\usepackage{multirow}
\usepackage[T1]{fontenc}
\usepackage[utf8]{inputenc}
\usepackage{lmodern}

% ── Hyperref (load last) ─────────────────────────────────────
\usepackage[
  colorlinks=true,
  linkcolor=blue!60!black,
  citecolor=green!50!black,
  urlcolor=blue!70!black,
  pdftitle={Context-Aware Vocabulary Pruning for LLM Inference Acceleration},
  pdfauthor={}
]{hyperref}

% ── Caption style ────────────────────────────────────────────
\captionsetup{font=small, labelfont=bf, skip=4pt}

% ── Section spacing ──────────────────────────────────────────
\usepackage{titlesec}
\titlespacing*{\section}{0pt}{8pt plus 2pt minus 1pt}{4pt plus 1pt}
\titlespacing*{\subsection}{0pt}{6pt plus 1pt}{3pt plus 1pt}
\titlespacing*{\subsubsection}{0pt}{4pt plus 1pt}{2pt plus 1pt}

% ── Paragraph spacing ────────────────────────────────────────
\setlength{\parskip}{2pt}
\setlength{\parindent}{1em}

% ── Abstract ─────────────────────────────────────────────────
\usepackage{abstract}
\renewcommand{\abstractnamefont}{\normalfont\bfseries\small}
\renewcommand{\abstracttextfont}{\small}
\setlength{\absleftindent}{0pt}
\setlength{\absrightindent}{0pt}

% ── Code / monospace ─────────────────────────────────────────
\usepackage{listings}
\lstset{basicstyle=\ttfamily\footnotesize, breaklines=true}

% ── Inline code helper ───────────────────────────────────────
\newcommand{\code}[1]{\texttt{#1}}

% ── Table helpers ────────────────────────────────────────────
\newcommand{\tblhead}[1]{\textbf{#1}}
\setlength{\tabcolsep}{4pt}

% ── Title block ──────────────────────────────────────────────
\title{%
  \vspace{-4mm}%
  \LARGE\bfseries
  Context-Aware Vocabulary Pruning for\\[2pt]
  LLM Inference Acceleration:\\[2pt]
  A Dual-Encoder Retrieval Approach%
  \vspace{-2mm}
}
\author{%
  \normalsize
  GPU benchmarks complete --- March 2026%
}
\date{}

% ============================================================
\begin{document}
% ============================================================

\maketitle
\thispagestyle{empty}

% ── Abstract ─────────────────────────────────────────────────
\begin{abstract}
Autoregressive decoding in large language models is dominated by two costs at
each token step: KV-cache attention and the \code{lm\_head} vocabulary
projection --- a matrix multiply that reads the full $[V \times d]$ weight
matrix from GPU memory on every step. At batch size 1, this operation is
\textbf{memory-bandwidth-bound}: the GPU's arithmetic units sit idle while
525\,MB of weights are streamed from VRAM per token. We address this
bottleneck through \textbf{dynamic vocabulary pruning}, restricting the
\code{lm\_head} projection to a context-dependent shortlist of $K$ candidate
tokens at each step, reducing memory traffic by up to 99.6\%.

We propose a \textbf{Dual-Encoder Router} --- a lightweight MLP (0.06\% of
model parameters) trained with Multiple Negatives Ranking Loss to map the
model's hidden state to a query vector, then retrieve the top-$K$ token
candidates by nearest-neighbour search over a pre-indexed corpus of sentence
completions. We benchmark four routing strategies against a full-vocabulary
baseline on WikiText-2 (\code{Llama-3.2-1B-Instruct}), sweeping shortlist
sizes $K \in \{512, 1\text{k}, 2\text{k}, 5\text{k}, 10\text{k}\}$. On an
NVIDIA T4 GPU the \code{lm\_head} accounts for 6.1\% of per-token latency;
the primary demonstrated gain is a 92--99.6\% reduction in \code{lm\_head}
memory bandwidth. Wall-clock speedup requires GPU-native kernels or larger
models (70B+) where the memory-bound regime is more pronounced.
\end{abstract}

% ────────────────────────────────────────────────────────────
\section{Introduction}
% ────────────────────────────────────────────────────────────

\subsection{The Autoregressive Bottleneck}

Modern LLMs generate text one token at a time. Despite the fact that the
transformer's attention mechanism is highly parallelisable over the
\emph{input} sequence, the \emph{generation} process is fundamentally
sequential: each new token depends on all previously generated tokens.

Within each step of this loop, the two most expensive operations are:
\begin{enumerate}[leftmargin=*, itemsep=0pt, topsep=2pt]
  \item \textbf{KV-cache attention} --- computing attention over the growing
        sequence of past tokens.
  \item \textbf{Vocabulary projection (\code{lm\_head})} --- projecting the
        final hidden state onto the full vocabulary to obtain next-token
        logits.
\end{enumerate}

On CPU and Apple Silicon (MPS) --- where memory bandwidth rather than raw
FLOP throughput is the binding constraint --- the \code{lm\_head} matrix
multiplication can account for a disproportionate share of per-token latency.
At batch size~1 --- the regime of interactive and agentic inference --- the
\code{lm\_head} projection is memory-bandwidth-bound: the arithmetic intensity
is too low to keep the GPU's compute units saturated, so the bottleneck is the
rate at which the weight matrix can be read from VRAM.

\subsection{The Wasted Computation Insight}

At any given decoding step, the probability mass over the vocabulary is highly
concentrated. In predictable contexts (e.g., ``The Eiffel Tower is located
in\ldots'') the top-1 token may carry $>$99\% probability. The model still
computes logits for all 32k--128k vocabulary entries. This motivates a natural
question: \emph{can the model's own internal representation predict which
vocabulary region is relevant before the full projection is run?}

\subsection{Contributions}

\begin{enumerate}[leftmargin=*, itemsep=0pt, topsep=2pt]
  \item \textbf{Profiling.} We quantify the \code{lm\_head} fraction of
        per-token latency for a 1B open-source LLM on an NVIDIA T4 GPU.
  \item \textbf{Router taxonomy.} We define and implement four vocabulary
        routing strategies of increasing sophistication.
  \item \textbf{Dual-Encoder Router.} A lightweight MLP trained with MNRL
        that maps the LLM's hidden state to a compact token shortlist,
        operating entirely within the model's own latent space.
  \item \textbf{Benchmark.} All methods evaluated on WikiText-2, reporting
        token acceptance rate, perplexity, and tokens/sec across
        $K \in \{512,1\text{k},2\text{k},5\text{k},10\text{k}\}$.
\end{enumerate}

% ────────────────────────────────────────────────────────────
\section{Background}
% ────────────────────────────────────────────────────────────

\subsection{Self-Attention and the Hidden State}

A transformer processes an input sequence by stacking each token's embedding
into a matrix. At each layer, self-attention allows every token's
representation to incorporate information from all preceding tokens. By the
final layer, the last row of the hidden state matrix --- the \textbf{``vibe''} ---
summarises the full context:
\[
  h_T = H^{(L)}_{T,:} \;\in\; \mathbb{R}^{d}.
\]

\subsection{The \texttt{lm\_head} as a Bottleneck}

The language modelling head maps the hidden state to a distribution over the
vocabulary:
\begin{equation}
  \text{logits} = h_T W_V^\top \;\in\; \mathbb{R}^{|V|}
\end{equation}
where $W_V \in \mathbb{R}^{|V| \times d}$ is the vocabulary embedding matrix.
At 16-bit precision, $W_V$ for Llama-3.2-1B occupies
$128{,}256 \times 2{,}048 \times 2\,\text{bytes} = \mathbf{525\,\text{MB}}$.
Reading 525\,MB per token step creates a hard ceiling on achievable
tokens/sec independent of model depth.

\subsection{Speculative Decoding and Related Work}

\textbf{Speculative decoding}~\cite{leviathan2023,chen2023} accelerates
autoregressive generation by interposing a small draft model that proposes
candidate tokens. A large verifier model evaluates the entire candidate
sequence in a single parallel forward pass.

Extensions include \textbf{Medusa}~\cite{cai2024}, which attaches multiple
decoding heads to the base model, and \textbf{tree-based speculative
decoding}, which generates a tree of candidate sequences from multiple draft
models. Our approach is orthogonal to all of these: it targets the
\code{lm\_head} projection cost directly and can be combined with any of the
above methods.

\subsection{Dual-Encoders and Contrastive Learning}

A \textbf{Dual-Encoder} processes two inputs through separate encoder towers
and learns to embed related pairs close together in a shared vector space.
\textbf{Multiple Negatives Ranking Loss (MNRL)}~\cite{henderson2017,karpukhin2020}
is the standard training objective. Given a batch of $(a_i, p_i)$
anchor-positive pairs:
\begin{equation}
  \mathcal{L} = -\frac{1}{B}\sum_{i=1}^{B}
    \log\frac{\exp(\text{sim}(a_i,p_i)/\tau)}
             {\sum_{j=1}^{B}\exp(\text{sim}(a_i,p_j)/\tau)}
\end{equation}
where $\text{sim}(\cdot,\cdot)$ is cosine similarity and $\tau$ is a
temperature parameter. All other positives in the batch serve as in-batch
negatives, requiring no explicit negative mining.

% ────────────────────────────────────────────────────────────
\section{Method}
% ────────────────────────────────────────────────────────────

\subsection{Problem Formalisation}

Let $f_\theta$ be a frozen autoregressive LLM with vocabulary $V$, model
dimension $d$, and vocabulary projection matrix $W_V$. At decoding step $t$:
\[
  \hat{y}_t = \arg\max_{v \in V}(h_t W_V^\top)_v.
\]
\textbf{Goal}: learn a router $r_\phi : \mathbb{R}^d \to 2^V$ such that
(i) $\hat{y}_t \in r_\phi(h_t)$ with high probability (\emph{acceptance
rate}), (ii) $|r_\phi(h_t)| \ll |V|$ (\emph{compression ratio}), and
(iii) the latency of computing $r_\phi(h_t)$ is small relative to the savings.

\subsection{Router Taxonomy}

We implement four routing strategies of increasing sophistication
(Table~\ref{tab:routers}).

\begin{table}[h]
\centering
\caption{Router taxonomy.}
\label{tab:routers}
\small
\begin{tabular}{@{}lp{4.0cm}c@{}}
\toprule
\tblhead{Router} & \tblhead{Description} & \tblhead{Training} \\
\midrule
Static Top-K & Highest $\|W_V[v,:]\|_2$ (context-free) & None \\
Cosine       & Highest $\cos(h_t, W_V[v,:])$           & None \\
Cluster      & Pre-cluster $W_V$; top-$C$ clusters by $\cos(h_t,\mu_c)$ & Offline \\
Dual-Encoder & Learned MLP; ANN search over $W_V$      & Yes \\
\bottomrule
\end{tabular}
\end{table}

\subsection{Dual-Encoder Router Design}

\subsubsection{Architecture}
A lightweight MLP projection head $g_\phi : \mathbb{R}^d \to \mathbb{R}^{d_r}$
is trained to map the frozen LLM's hidden state to a query vector in a shared
embedding space. The vocabulary token embeddings $W_V$ serve as the corpus
searched at inference time. Both query and keys live in the LLM's own latent
space; no cross-model alignment is required.

\subsubsection{Training Data Construction}
From the WikiText-2 training split: (1) tokenise each sentence; (2) for each
sentence of length $T \geq 8$, sample a split point
$s \sim \text{Uniform}(\lfloor 0.4T\rfloor, \lfloor 0.7T\rfloor)$;
(3) \emph{anchor}: run the frozen LLM on the prefix and extract $h_s$;
(4) \emph{positive}: the ground-truth next token $y_{s+1}$, represented as
$W_V[y_{s+1},:]$.

\subsubsection{Training Objective}
\begin{equation}
  \mathcal{L} = -\frac{1}{B}\sum_{i=1}^{B}
    \log\frac{\exp(\cos(g_\phi(h_i),W_V[y_i,:])/\tau)}
             {\sum_{j=1}^{B}\exp(\cos(g_\phi(h_i),W_V[y_j,:])/\tau)}
\end{equation}
with in-batch negatives (all $y_j \neq y_i$ in the batch).

\subsubsection{Inference Integration}
(1) Pre-compute and cache all rows of $W_V$ normalised to unit length.
(2) At step $t$: compute $q_t = g_\phi(h_t)$; retrieve the $K$ token indices
with highest $\cos(q_t, W_V[v,:])$ via brute-force matmul.
(3) Slice $W_V$ to the shortlist and compute logits only over those $K$ rows.

% ────────────────────────────────────────────────────────────
\section{Experimental Setup}
% ────────────────────────────────────────────────────────────

\subsection{Model}
\code{meta-llama/Llama-3.2-1B-Instruct} --- 1B parameters, 2048-dim hidden
state, 128,256-token vocabulary, 16 transformer layers, \code{bfloat16}.

\subsection{Dataset}
\textbf{WikiText-2-raw-v1}~\cite{merity2016}: standard language modelling
benchmark; plain English prose. Router training: \code{train} split
(${\sim}2$M tokens). Perplexity evaluation: \code{test} split
(${\sim}240$k tokens). Latency benchmark: 50 randomly sampled contiguous
passages of 128 tokens from the \code{validation} split.

\subsection{Hardware}
NVIDIA T4 GPU (16\,GB VRAM), CUDA backend, \code{bfloat16}. All latency
measurements taken after a 5-step warm-up; reported as mean over 50 prompts.

\subsection{Evaluation Metrics}

\begin{itemize}[leftmargin=*, itemsep=1pt, topsep=2pt]
  \item \textbf{Token acceptance rate}: fraction of steps where the
        ground-truth token is in the shortlist.
  \item \textbf{Perplexity (PPL)}: $\exp\!\bigl(-\frac{1}{N}\sum_t \log
        P_\theta(y_t \mid y_{<t})\bigr)$ over the \code{test} split.
  \item \textbf{$\Delta$PPL}: $\text{PPL}_{\text{pruned}} -
        \text{PPL}_{\text{baseline}}$.
  \item \textbf{Tokens/sec}: wall-clock generation speed, mean over 50 prompts.
  \item \textbf{\code{lm\_head} time \%}: fraction of per-token time spent in
        the vocabulary projection.
\end{itemize}

% ────────────────────────────────────────────────────────────
\section{Results}
% ────────────────────────────────────────────────────────────

All experiments: NVIDIA T4 GPU, CUDA, \code{bfloat16},
\code{Llama-3.2-1B-Instruct}. Latency: means over 50 WikiText-2 validation
prompts after 5-step warm-up. Perplexity: full WikiText-2 test split.

\subsection{Baseline Profiling}

\begin{table}[h]
\centering
\caption{Full-vocabulary baseline measurements (NVIDIA T4 GPU).}
\label{tab:baseline}
\small
\begin{tabular}{@{}lr@{}}
\toprule
\tblhead{Metric} & \tblhead{Value} \\
\midrule
Tokens/sec (full vocab)        & \textbf{26.60} \\
Mean step time                 & 37.60\,ms \\
\code{lm\_head} time (mean)    & 2.30\,ms \\
\code{lm\_head} time \%        & \textbf{6.1\%} \\
Baseline perplexity (WikiText-2 test) & \textbf{26.24} \\
\code{lm\_head} weight matrix  & 525\,MB \\
\bottomrule
\end{tabular}
\end{table}

The \code{lm\_head} accounts for only 6.1\% of per-token latency. This is a
critical constraint: any router whose overhead exceeds 2.30\,ms will make
pruning net-negative on a 1B model (Section~\ref{sec:lmhead-fraction}).

\begin{figure}[H]
\centering
\includegraphics[width=\columnwidth]{../results/baseline_profile.png}
\caption{Baseline decode step profiling: per-operation time breakdown.}
\label{fig:baseline}
\end{figure}

\subsection{Token Acceptance Rate vs.\ K}

Acceptance rate is the fraction of decoding steps where the ground-truth next
token falls within the pruned shortlist (Table~\ref{tab:acceptance}).

\begin{table*}[t]
\centering
\caption{Token acceptance rate (\%) vs.\ shortlist size $K$ for all routers.}
\label{tab:acceptance}
\small
\begin{tabular}{@{}lrrrrr@{}}
\toprule
\tblhead{Router} & \tblhead{K=512} & \tblhead{K=1k} & \tblhead{K=2k} & \tblhead{K=5k} & \tblhead{K=10k} \\
\midrule
Static Top-K         & 0.09 & 0.25 & 0.50 & 2.5  & 7.7  \\
Cluster              & 33.8 & 54.8 & 68.1 & 76.6 & 83.3 \\
Cosine               & 98.4 & 98.4 & 98.4 & 98.4 & 98.4 \\
Dual-Encoder (step)  & 99.5 & 99.5 & 99.5 & 99.5 & 99.5 \\
MLP Graph            & 99.9 & 100.0 & 100.0 & 100.0 & 100.0 \\
Attention Graph      & 82.7 & 99.8 & 99.8 & 99.8 & 99.8 \\
\bottomrule
\end{tabular}
\end{table*}

Key observations: Static Top-K is near-random ($<$8\% at K=10k). Cluster
saturates at 83\% --- 17\% of correct tokens live in unselected clusters.
Cosine acceptance is flat at 98.4\% across all $K$ values, indicating the
router reaches its information-theoretic ceiling early. Dual-Encoder (step)
and MLP Graph achieve the highest acceptance (99.5\% and $\geq$99.9\%
respectively), essentially independent of $K$ above 512.

\begin{figure*}[t]
\centering
\includegraphics[width=\textwidth]{../results/acceptance_rate_vs_k.png}
\caption{Token acceptance rate vs.\ shortlist size $K$ for all routers.}
\label{fig:acceptance}
\end{figure*}

\subsection{Perplexity Degradation vs.\ K}

$\Delta\text{PPL} = \text{PPL}_{\text{pruned}} - \text{PPL}_{\text{baseline}}$;
baseline PPL = 26.24 (Table~\ref{tab:ppl}).

\begin{table*}[t]
\centering
\caption{Perplexity degradation ($\Delta$PPL) vs.\ shortlist size $K$.
         Static Top-K and Cluster omitted due to extreme values
         ($\Delta$PPL $> 10^6$).}
\label{tab:ppl}
\small
\begin{tabular}{@{}lrrrrr@{}}
\toprule
\tblhead{Router} & \tblhead{K=512} & \tblhead{K=1k} & \tblhead{K=2k} & \tblhead{K=5k} & \tblhead{K=10k} \\
\midrule
Cosine                          & +19.00 & +10.83 & +6.79 & +3.54 & \textbf{+2.73} \\
Dual-Encoder (step)             & +24.64 & +8.95  & +3.93 & +1.68 & \textbf{+1.03} \\
Dual-Encoder (prefetch+refresh)\textsuperscript{$\dagger$} & --- & --- & --- & +146.15 & --- \\
MLP Graph                       & +39.59 & +13.22 & +6.39 & +2.89 & \textbf{+2.11} \\
Attention Graph                 & +15025 & +21.58 & +5.91 & +2.10 & \textbf{+1.33} \\
\bottomrule
\multicolumn{6}{l}{\footnotesize $\dagger$ Prefetch+refresh uses a fixed
  effective $K{\approx}4{,}634$ averaged over the evaluation.}
\end{tabular}
\end{table*}

At K=10k the three best methods cluster closely: DE step (+1.03), Attention
Graph (+1.33), MLP Graph (+2.11), and Cosine (+2.73) --- all within 11\% of
baseline PPL. DE step at K=5k achieves $\Delta$PPL=+1.68 with a 96.1\%
memory bandwidth reduction.

\begin{figure*}[t]
\centering
\includegraphics[width=\textwidth]{../results/ppl_degradation_vs_k.png}
\caption{Perplexity degradation ($\Delta$PPL) vs.\ shortlist size $K$.
         Static Top-K and Cluster are omitted due to extreme values.}
\label{fig:ppl}
\end{figure*}

\subsection{Tokens/sec vs.\ K}

Wall-clock throughput on an NVIDIA T4 GPU. Baseline = 26.60\,tok/s
(Table~\ref{tab:throughput}).

\begin{table*}[t]
\centering
\caption{Wall-clock throughput (tok/s) vs.\ shortlist size $K$.
         Baseline = 26.60\,tok/s (NVIDIA T4 GPU).}
\label{tab:throughput}
\small
\begin{tabular}{@{}lrrrrrr@{}}
\toprule
\tblhead{Router} & \tblhead{K=512} & \tblhead{K=1k} & \tblhead{K=2k} & \tblhead{K=5k} & \tblhead{K=10k} & \tblhead{vs.\ baseline} \\
\midrule
\textbf{Baseline} & --- & --- & --- & --- & 26.60 & --- \\
Static Top-K      & 26.07 & 26.40 & 26.36 & 26.82 & 26.70 & $-1\%$ to $+1\%$ \\
Cluster           & 24.34 & 24.18 & 24.25 & 24.27 & 23.99 & $-9\%$ to $-10\%$ \\
Cosine            & 24.27 & 24.17 & 23.77 & 23.51 & 23.72 & $-9\%$ to $-12\%$ \\
DE (prefetch+refresh) & --- & --- & --- & \textbf{27.67} & --- & $\mathbf{+4\%}$ \\
Dual-Encoder (step)   & 23.27 & 22.97 & 23.20 & 23.16 & 22.39 & $-13\%$ to $-16\%$ \\
MLP Graph         & 24.28 & 24.31 & 23.81 & 22.89 & 22.83 & $-14\%$ to $-16\%$ \\
Attention Graph   & 22.72 & 20.71 & 21.30 & 20.12 & 19.98 & $-15\%$ to $-25\%$ \\
\bottomrule
\end{tabular}
\end{table*}

On GPU the throughput gap between pruned methods and baseline is much narrower
than on memory-limited hardware. Static Top-K adds negligible overhead (no
neural network pass), sitting within $\pm$1\% of baseline. The DE
(prefetch+refresh) variant is the only pruned method to exceed baseline
throughput (+4\%, 27.67\,tok/s), by amortising router cost across many
decoding steps --- though at a severe quality penalty (+146 PPL). Attention
Graph remains the worst due to \code{output\_attentions=True} forcing eager
attention, but the penalty ($-15\%$ to $-25\%$) is far smaller than on
memory-bandwidth-limited hardware.

\begin{figure*}[t]
\centering
\includegraphics[width=\textwidth]{../results/throughput_vs_k.png}
\caption{Wall-clock tokens/sec vs.\ shortlist size $K$.
         Dashed line = full-vocab baseline (26.60\,tok/s, NVIDIA T4 GPU).}
\label{fig:throughput}
\end{figure*}

\subsection{Pareto Frontier: Quality vs.\ Speed}

The accuracy--throughput Pareto frontier identifies configurations not
dominated --- no other configuration achieves both higher throughput and lower
PPL degradation simultaneously.

\begin{table}[h]
\centering
\caption{Approximate Pareto frontier among pruned methods (T4 GPU).}
\label{tab:pareto}
\small
\begin{tabular}{@{}llrrl@{}}
\toprule
\tblhead{Method} & \tblhead{K} & \tblhead{Tok/s} & \tblhead{$\Delta$PPL} & \tblhead{Note} \\
\midrule
DE (prefetch) & 4634 & 27.67 & $+146$ & Fastest; quality unusable \\
Cluster     & 10k   & 23.99 & $+834$ & Quality unusable \\
Cosine      & 10k  & 23.72 & $+2.73$ & Good quality \\
\textbf{DE (step)} & \textbf{5k} & \textbf{23.16} & \textbf{+1.68} & \textbf{Best trade-off} \\
DE (step)   & 10k  & 22.39 & $+1.03$ & Marginal PPL gain \\
\bottomrule
\end{tabular}
\end{table}

On T4, DE (step) at K=5k represents the recommended operating point:
23.16\,tok/s ($-$13\% vs.\ baseline) with $\Delta$PPL=+1.68 --- 96.1\%
memory bandwidth reduction at minimal quality cost. The throughput penalty
is considerably more acceptable on GPU than on memory-bandwidth-limited
hardware.

\begin{figure}[H]
\centering
\includegraphics[width=\columnwidth]{../results/pareto_frontier.png}
\caption{Pareto frontier: throughput vs.\ $\Delta$PPL. Points toward the
  top-left are preferred. Full-vocab baseline sits at the optimal top-left
  corner.}
\label{fig:pareto}
\end{figure}

\begin{figure}[H]
\centering
\includegraphics[width=\columnwidth]{../results/shortlist_growth.png}
\caption{Mean shortlist size growth vs.\ $K$ (DE prefetch+refresh variant).}
\label{fig:shortlist}
\end{figure}

\subsection{Graph-Based Routers}

\textbf{MLP Transition Graph.} A directed graph over vocabulary tokens where
edge weights encode how frequently token $v$ follows token $u$ in training.
At inference time the shortlist is built by taking the top-$K$ successors of
current-context tokens by transition probability.
\begin{itemize}[leftmargin=*, itemsep=0pt, topsep=2pt]
  \item Acceptance rate: $\geq$99.9\% across all $K$ --- the highest of any method.
  \item Throughput: 22.83--24.31\,tok/s ($-14\%$ to $-16\%$ vs.\ baseline).
  \item PPL at K=10k: +2.11.
\end{itemize}

\textbf{Attention Graph.} Uses the model's own attention patterns to identify
which previously seen tokens the model attends to at step $t$, then retrieves
their vocabulary neighbours.
\begin{itemize}[leftmargin=*, itemsep=0pt, topsep=2pt]
  \item Acceptance at $K \geq 1$k: 99.8\%.
  \item Throughput: 19.98--22.72\,tok/s ($-15\%$ to $-25\%$ vs.\ baseline).
        \code{output\_attentions=True} forces eager attention; on GPU this
        penalty is proportionally smaller than on memory-bandwidth-limited
        hardware.
  \item PPL at K=10k: +1.33 --- second-best quality result after DE step.
\end{itemize}

The attention graph's quality ($\Delta$PPL=+1.33) is the second best of all
methods. A CUDA-native implementation reading pre-computed attention patterns
without re-running attention would recover this overhead.

\subsection{Speculative Decoding via Retrieval}

Two speculative decoding approaches were implemented and benchmarked.

\subsubsection{Chain-based (v1)}
Token 1 selected from the combined Attn+Graph shortlist argmax; tokens 2..D
from a 1-hop graph walk. The full model runs two forward passes per step, and
\code{output\_attentions=True} forces eager attention throughout.

\begin{table}[h]
\centering
\caption{Chain-based speculative decoding results (T4 GPU).}
\label{tab:spec-v1}
\small
\begin{tabular}{@{}lrrr@{}}
\toprule
\tblhead{Config} & \tblhead{Tok/s} & \tblhead{Draft acc.} & \tblhead{Speedup} \\
\midrule
Cosine ($D{=}4$)          & 2.92  & 27.7\% & $0.11{\times}$ \\
Attn+Graph ($D{=}4$)      & 2.88  & 28.1\% & $0.11{\times}$ \\
\textbf{Full-vocab baseline} & \textbf{26.60} & --- & \textbf{1.0×} \\
\bottomrule
\end{tabular}
\end{table}

At 2.88--2.92\,tok/s this is ${\sim}$11\% of baseline. Two compounding
factors: (1) two full forward passes per step with no parallelism gain;
(2) \code{output\_attentions=True} adding eager-attention overhead per step.

\begin{figure}[H]
\centering
\includegraphics[width=\columnwidth]{../results/speculative_decoding_sweep.png}
\caption{Speculative decoding (chain-based v1): effective tok/s and draft
  acceptance rate vs.\ draft length $D$.}
\label{fig:spec-v1}
\end{figure}

\subsubsection{Retrieval-based (v2)}
An offline index of 50k completion sequences (8 tokens each) is built by
extracting $h_T$ at split points across WikiText-2 training windows. At
inference time: (1) one probe forward pass with KV-cache captures $h_T$;
(2) top-$K$ completions retrieved by cosine similarity; (3) each candidate
verified using the cached KV-cache --- $O(D)$ attention cost instead of
$O(T+D)$. No \code{output\_attentions}; standard SDPA throughout.

\begin{table}[h]
\centering
\caption{Retrieval-based speculative decoding (v2) sweep results (T4 GPU).}
\label{tab:spec-v2}
\small
\begin{tabular}{@{}rrrrrr@{}}
\toprule
\tblhead{K} & \tblhead{D} & \tblhead{Tok/s} & \tblhead{Tok/step} & \tblhead{Draft acc.} & \tblhead{Speedup} \\
\midrule
 4 &  4 & 2.44 & 1.19 & 1.2\% & $0.09{\times}$ \\
 4 &  8 & 2.38 & 1.16 & 0.5\% & $0.09{\times}$ \\
 4 & 16 & 2.38 & 1.16 & 0.3\% & $0.09{\times}$ \\
 8 &  4 & 1.63 & 1.27 & 0.8\% & $0.06{\times}$ \\
 8 &  8 & 1.58 & 1.24 & 0.4\% & $0.06{\times}$ \\
 8 & 16 & 1.58 & 1.24 & 0.2\% & $0.06{\times}$ \\
16 &  4 & 1.03 & 1.40 & 0.6\% & $0.04{\times}$ \\
16 &  8 & 0.97 & 1.34 & 0.3\% & $0.04{\times}$ \\
16 & 16 & 0.97 & 1.34 & 0.1\% & $0.04{\times}$ \\
32 &  4 & 0.63 & 1.59 & 0.5\% & $0.02{\times}$ \\
\bottomrule
\end{tabular}
\end{table}

The KV-cache fix eliminated the $O(T^2)$ quadratic cost from v1, improving
throughput to 0.63--2.44\,tok/s. However, draft acceptance rates collapsed to
\textbf{0.1--1.2\%} --- far below the ${\sim}$70--80\% needed for speculative
speedup. The root cause is a \textbf{retrieval signal mismatch}: the index
embeds prefix hidden states from an offline training-data pass, while at
inference time $h_T$ is a function of the full prompt history. The dot-product
similarity is insufficient to recover semantically valid continuations, so
nearly every draft token is rejected. The retrieval signal needs to be
\emph{discriminatively trained}, not borrowed from the LLM's general-purpose
representation.

\subsection{Cost Analysis}

Using GPU on-demand pricing (A100 @ \$3.00/hr = \$0.000833/s):

\begin{table}[h]
\centering
\caption{Cost per 1M tokens at A100 on-demand pricing (tok/s from T4 GPU).}
\label{tab:cost}
\small
\begin{tabular}{@{}llrrl@{}}
\toprule
\tblhead{Method} & \tblhead{K} & \tblhead{Tok/s} & \tblhead{\$/1M tok} & \tblhead{vs.\ base} \\
\midrule
\textbf{Baseline} & 128k & 26.60 & \$31.33 & --- \\
Cosine       & 10k & 23.72 & \$35.13  & $+12\%$ \\
DE (step)    &  2k & 23.20 & \$35.92  & $+15\%$ \\
DE (step)    &  5k & 23.16 & \$35.99  & $+15\%$ \\
DE (step)    & 10k & 22.39 & \$37.22  & $+19\%$ \\
MLP Graph    & 10k & 22.83 & \$36.50  & $+16\%$ \\
Attn Graph   & 10k & 19.98 & \$41.72  & $+33\%$ \\
Spec (v1)    & D=4 &  2.88 & $\sim$\$289 & $+822\%$ \\
\bottomrule
\end{tabular}
\end{table}

On GPU the cost overhead of pruned methods is much more modest than on
memory-limited hardware: 12--33\% above baseline for the viable methods.
Break-even between router cost and \code{lm\_head} savings still requires a
larger model (see Section~\ref{sec:lmhead-fraction}), but the arithmetic is
less unfavourable on GPU.

% ────────────────────────────────────────────────────────────
\section{Discussion}
% ────────────────────────────────────────────────────────────

\subsection{The Blind-Spot Risk}

Dynamic vocabulary pruning introduces a fundamental tension: the router can
only shortlist tokens it predicts are relevant, but creative or unexpected
tokens by definition have low prior probability given context. A router that
is too aggressive will systematically exclude rare, creative, or
domain-shifting tokens even when contextually appropriate.

Our results directly quantify this risk. DE (step) achieves 99.5\% acceptance
at all $K$ values from 512 to 10k --- meaning the correct token is in the
shortlist 99.5\% of steps regardless of $K$. The limiting factor is not
shortlist size but which embedding-space neighbours the router retrieves: the
0.5\% of missed tokens are systematically hard cases (rare tokens, domain
shifts) that the router misses at any $K$.

\subsection{Limitations}

\begin{itemize}[leftmargin=*, itemsep=1pt, topsep=2pt]
  \item Evaluation limited to a single 1B model and a single dataset
        (WikiText-2 English prose).
  \item Evaluation on a single GPU tier (NVIDIA T4); results at
        higher-bandwidth GPUs (A100, H100) may differ.
  \item The dual-encoder is trained on token-level prediction; it may not
        generalise to multi-token idiomatic phrases or specialised domains
        (code, mathematics, non-English text).
  \item Router training uses only WikiText-2; domain mismatch could reduce
        acceptance rates.
\end{itemize}

\subsection{The \texttt{lm\_head} Fraction Problem}
\label{sec:lmhead-fraction}

The core finding --- that vocabulary pruning does not improve throughput on a
1B model --- follows directly from arithmetic:
\begin{align*}
  \text{lm\_head saving} &= 37.60\,\text{ms} \times 0.061 \times 0.922 = 2.12\,\text{ms} \\
  \text{Router overhead (DE, K=10k)} &= 44.67 - 37.60 = 7.07\,\text{ms} \\
  \text{Net:} &\;{+}4.95\,\text{ms/step} \;\Rightarrow\; {-}13\%\;\text{throughput}
\end{align*}
For pruning to be net-positive, the \code{lm\_head} fraction must exceed:
\begin{equation}
  f_{\text{break-even}} =
  \frac{\text{router\_overhead\_ms}}{\text{step\_ms} \times (1 - K/V)}
\end{equation}
On this 1B model, break-even requires $f \approx 20\%$; measured value is
6.1\%. At 70B scale (step $\approx$1,200\,ms), break-even drops to
${\sim}0.64\%$ --- well below the expected \code{lm\_head} fraction at any
model size. This is the regime where vocabulary pruning becomes genuinely
useful.

Importantly, the \textbf{memory bandwidth saving} is real at all scales.
Reducing \code{lm\_head} traffic by 92\% matters for edge deployment (models
that barely fit in VRAM) and for batched inference where memory bandwidth is
the binding resource.

\subsection{Speculative Decoding via Retrieval}

Two implementations were evaluated. The chain-based v1 achieved 28\% draft
acceptance while penalised by two full forward passes per step and eager
attention overhead. The retrieval-based v2 eliminated both overhead sources
via KV-cache reuse and standard SDPA, but draft acceptance collapsed to
0.2--1.1\%. Key factors:

\begin{enumerate}[leftmargin=*, itemsep=1pt, topsep=2pt]
  \item \textbf{Retrieval signal mismatch.} Index embeddings are from
        offline training-data prefixes; inference-time $h_T$ differs
        systematically. Cosine similarity is low-signal for valid continuations.
  \item \textbf{No discriminative training.} The index uses general-purpose
        LLM representations, not representations trained to distinguish
        accepted from rejected draft sequences.
  \item \textbf{Sequential verifier passes.} A tree-structured verifier pass
        over all $K$ candidates in a single batched forward pass (Option C)
        would reduce per-step cost from $O(K \cdot D)$ to $O(D)$.
\end{enumerate}

The underlying architecture --- retrieval of completion sequences, KV-cache
reuse in the verifier, standard SDPA --- is sound. The missing ingredient is a
discriminatively trained retrieval model.

% ────────────────────────────────────────────────────────────
\section{Conclusion and Future Work}
% ────────────────────────────────────────────────────────────

We implemented and benchmarked six vocabulary pruning strategies and two
retrieval-based speculative decoders against a full-vocabulary baseline on
\code{Llama-3.2-1B-Instruct} and WikiText-2, on an NVIDIA T4 GPU.

\textbf{The memory bandwidth result is the genuine contribution.} Pruning
\code{lm\_head} to K=10k reduces the weight matrix loaded per step from
525\,MB to 41\,MB --- a 92.2\% reduction --- while DE (step) at K=5k keeps
$\Delta$PPL to +1.68 above the 26.24 baseline.

\textbf{The throughput result is an honest negative.} On a 1B model on a T4
GPU, the \code{lm\_head} fraction is only 6.1\%. Break-even requires
$f \approx 20\%$ vs.\ measured 6.1\%; no pruned method reduces wall-clock
cost at this model scale. However, the overhead gap is considerably narrower
than on memory-bandwidth-limited hardware: DE (step) at K=5k costs only
$-$13\% throughput vs.\ baseline, and the DE (prefetch+refresh) variant
exceeds baseline throughput (+4\%, 27.67\,tok/s) by amortising router cost
across many steps, at the cost of +146 PPL degradation. Break-even for
quality-preserving pruning requires a much larger model (70B+) or a
GPU-native router kernel reducing overhead below ${\sim}$1\,ms.

\textbf{Speculative decoding via retrieval} was evaluated in two
implementations. The chain-based v1 achieved 2.88\,tok/s (11\% of baseline)
with 28\% draft acceptance. The retrieval-based v2 removed both overhead
sources but draft acceptance collapsed to 0.1--1.2\% due to retrieval signal
mismatch. A useful retrieval-based speculative decoder requires a
\emph{discriminatively trained} retrieval model.

\textbf{Future work:}
\begin{itemize}[leftmargin=*, itemsep=1pt, topsep=2pt]
  \item \textbf{GPU-native router kernel.} Implement the DE router as a fused
        CUDA kernel to reduce overhead from ${\sim}$7\,ms to $<$1\,ms --- the
        key step for making pruning viable on 7B models.
  \item \textbf{INSTRUCTOR-style router.} Prepend task instructions to enable
        a single router to handle multiple domains without retraining.
  \item \textbf{Discriminative retrieval for speculative decoding.} Train a
        contrastive retrieval model on accepted vs.\ rejected draft sequences.
        Combined with a compact Medusa head for verification, this is the most
        promising path to retrieval-based speculative speedup.
  \item \textbf{Online router adaptation.} Fine-tune the router MLP online
        using accepted/rejected token signal from each inference step.
  \item \textbf{Larger model evaluation.} Test 8B and 70B Llama-3.1
        checkpoints, where the \code{lm\_head} fraction is expected to be
        higher and router overhead is proportionally smaller.
  \item \textbf{Higher-bandwidth GPU evaluation.} Reproduce on A100/H100 to
        determine whether higher memory bandwidth changes the break-even
        arithmetic for smaller models.
\end{itemize}

% ────────────────────────────────────────────────────────────
\bibliographystyle{plain}
\begin{thebibliography}{9}

\bibitem{cai2024}
T.~Cai et al.
\newblock Medusa: Simple LLM inference acceleration framework with multiple
  decoding heads.
\newblock 2024.

\bibitem{chen2023}
C.~Chen et al.
\newblock Accelerating large language model decoding with speculative sampling.
\newblock 2023.

\bibitem{henderson2017}
M.~Henderson et al.
\newblock Efficient natural language response suggestion for smart reply.
\newblock 2017.

\bibitem{karpukhin2020}
V.~Karpukhin et al.
\newblock Dense passage retrieval for open-domain question answering.
\newblock In \emph{EMNLP}, 2020.

\bibitem{leviathan2023}
Y.~Leviathan, M.~Kalman, and Y.~Matias.
\newblock Fast inference from transformers via speculative decoding.
\newblock In \emph{ICML}, 2023.

\bibitem{merity2016}
S.~Merity et al.
\newblock Pointer sentinel mixture models.
\newblock 2016. (WikiText-2 dataset)

\bibitem{su2022}
H.~Su et al.
\newblock One embedder, any task: Instruction-finetuned text embeddings.
\newblock 2022.

\end{thebibliography}

\end{document}
